% P10-4a, 14.09.2026: Mechanismus, Quellenreichweite und v4-Nachweis angeglichen.
% M08, 11.09.2026: Objektbedeutung/F1 ergänzt; Status und Auflösung konzentriert.
% Gegenprüfung: F1 07-github/applied/core-before.json; DecisionResolution*
% sowie CoreKangal. Drei Auflösungsarten nur mit bestehendem Zielelement.
\section{Core-zentrierte Projektfortschreibung}
\label{sec:core-projektfortschreibung}

Der Core führt den maßgeblichen fachlichen Projektzustand über
Verarbeitungsläufe hinweg und konkretisiert damit die Anforderungen
A5 und A7. Seine Elemente besitzen fachliche Typen, Kennungen,
Beziehungen und eine Änderungshistorie. Requirements halten die
Anforderungen fest; Features bündeln ein Thema; Arbeitspakete
(PBIs) beschreiben die dafür geplante Arbeit. Architektur- und
Entscheidungselemente ergänzen Rahmenbedingungen und offene
Klärungsgegenstände. Daneben speichert das System begleitende
Governance-Informationen, etwa wartende oder abgelehnte Vorschläge.
Ihre Speicherung verleiht ihnen noch keine fachliche Geltung.

Im Besuchsfall F1 referenzieren die Anforderungen REQ-82 und REQ-83
denselben Ledger-Claim. Nach der freigegebenen Arbeitspaketbildung
gehören beide zum Feature FC-15; PBI-046 plant die Besuchsankündigung,
PBI-047 die Übersicht für Pflegende. Die jeweilige Beziehung
\emph{deckt} ordnet das Arbeitspaket seiner Anforderung zu.
Sie belegt geplante Zuständigkeit, keine bereits implementierte
Softwarefunktion und keine gemessene semantische Abdeckung.
Den gesamten Fall führt Tabelle~\ref{t:beispiel-besuchsplanung} zusammen.

Quellenangaben an einem Element und die Herkunft einer Beziehung
beantworten unterschiedliche Fragen. Die Quellenangaben verweisen
auf die zugrunde gelegten Eingaben. Eine Erzeugerangabe an der Beziehung
hält dagegen fest, welcher Verarbeitungsschritt die Verbindung
angelegt hat. Sie bezeichnet den Mechanismus, keinen Zeitpunkt;
für die zeitliche Rekonstruktion werden zusätzlich die archivierten
Laufartefakte benötigt. Tabelle~\ref{t:core-objektmodell} erläutert
die fachliche Bedeutung der Beziehungen und unterscheidet sie von
Quellenbezügen und der Zuordnung zu externen Issues.

% P10-4a, 14.09.2026: Mechanismus, Quellenreichweite und v4-Nachweis angeglichen.
% ============================================================
% Tabelle 5.5-1: Zentrale Objekte und Beziehungen des Core — v01 (09.09.)
% In Overleaf ablegen als: tables/core-objektmodell.tex
% Aufruf im Fließtext: \input{\tabledir/core-objektmodell.tex}
% Label: t:core-objektmodell
% Stil: chap6.1-Konvention.
% HERKUNFT: Relationstypen und Feldnamen aus state/core/
% project-state.json (relations: fromId/relationType/toId; items:
% sourceRunId/sourceClaimIds) — am Besuchsplanungs-Bestand verifiziert
% (REQ-82/83 · PBI-046/047 · FC-15 · gh#44/45; supersedes am
% E47-02-Fall). Bewusst NUR die tragenden Beziehungen, kein
% vollständiges Klassen-/Feldinventar (Umsetzungsplan P4).
% ============================================================

\begin{table}[htbp]
  \centering
  \small
  \begin{tabular}{>{\raggedright\arraybackslash}p{0.20\textwidth}>{\raggedright\arraybackslash}p{0.17\textwidth}>{\raggedright\arraybackslash}p{0.17\textwidth}>{\raggedright\arraybackslash}p{0.34\textwidth}}
    \hline
    \textbf{Ausgangsobjekt} & \textbf{Beziehung} & \textbf{Zielobjekt} & \textbf{Bedeutung} \\
    \hline
    Core-Element & Herkunftsangaben, je Eingang & Ledger-Claim oder Eingangs-/Übernahmebeleg & deklarierter Quellenzugang; Auflösbarkeit und Bezug zur aktuellen Fassung gesondert zu prüfen \\
    Arbeitselement (PBI) & deckt & Requirement / Architektur-Element & ordnet geplante Arbeit einem fachlichen Inhalt zu; kein Umsetzungs- oder Abdeckungsnachweis \\
    Requirement / PBI & gehört zu Feature & Feature & fachliche Themenstruktur \\
    Arbeitselement (PBI) & begrenzt durch & Architektur-Element & technische Rahmenbedingung wirkt auf die Umsetzung \\
    neue Fassung & löst ab & frühere Fassung & Ablösung mit erhaltener Historie (keine Löschung) \\
    Arbeitselement (PBI) & umgesetzt durch Issue & externes Issue & operative Darstellung des Arbeitspakets; keine implementierte Softwarefunktion \\
    \hline
  \end{tabular}
  \caption[Objekte und Beziehungen im Core]{Zentrale Objekte und
    Beziehungen im Core. Herkunfts- und Issue-Zuordnung verweisen
    nach außen; die übrigen Zeilen verbinden Core-Objekte.
    Begleitende Governance-Informationen sind nicht dargestellt.}
  \label{t:core-objektmodell}
\end{table}


% P10-6: Erzeugung von Elementen und Beziehungen als zusammenhängenden Absatz setzen.
\begin{samepage}
\emph{Vom Eingang zum verknüpften Bestand.} Der Einordnungsagent
erhält Werkzeuge, um eingehende Inhalte, vorhandene Core-Elemente,
offene Entscheidungen und frühere Ablehnungen nachzuschlagen.
Er wählt Abfragen und schlägt vor, ob eine Aussage neu ist,
Bekanntes bestätigt, verändert oder Klärung erfordert. Die
Suchergebnisse liefern Anhaltspunkte; die fachliche Zuordnung bleibt
ein Modellurteil. Sein Operationsplan verknüpft jeden Eingang mit
der vorgeschlagenen Wirkung und bei Änderungen mit einer bestehenden
Core-Kennung. Nach Prüfung und menschlicher Freigabe setzt die
deterministische Anwendung die ausgewählten Operationen um: Sie
vergibt neue Kennungen oder aktualisiert bestehende Elemente und
übernimmt die vorgesehenen Herkunftsangaben aus dem Eingang.
\par
\end{samepage}

Die fachliche Vernetzung wird in weiteren, eigens freigegebenen
Schritten ergänzt. Bei der Arbeitspaketbildung schlagen Modelle
Themenzuordnungen und den Zuschnitt der Arbeit vor. Die Anwendung
erzeugt daraus PBI-zu-Anforderung-Beziehungen und die
Feature-Zuordnungen. Wird später ein bestehendes PBI um eine
Anforderung erweitert, ergänzt sie dessen Beziehung \emph{deckt}
und ordnet die Anforderung dem Feature des PBI zu. Eine zusätzliche
freigegebene Angleichung kann Titel, Beschreibung und
Akzeptanzkriterien anpassen. Das Netz wächst somit durch konkrete
Übernahmeoperationen; eine neue Beziehung allein bedeutet noch
keine abgeschlossene Inhaltsangleichung.

\emph{Kennungen und Herkunft.} Evidence-Unit-Kennungen entstehen
positionsbasiert: Bei unveränderter Quelle und Segmentierung sind
sie wiederholbar, über Quelländerungen aber nicht stabil und zwischen
Quellen nicht eindeutig. Claim-Kennungen gelten innerhalb eines
Laufbestands. Derselbe Zeichenwert in einem neuen Lauf bezeichnet
deshalb nicht automatisch dieselbe Aussage. Laufübergreifende
fachliche Identität tragen die Core-Kennungen. Normalisierte
Textschlüssel unterstützen die Suche nach möglichen Entsprechungen;
die fachliche Zuordnung erfolgt im kontrollierten Einordnungs- und
Autorisierungsprozess. Eine garantierte Wiederzuordnung bei erneuter
Verarbeitung folgt daraus nicht. Beim Transkript verbindet der
Herkunftsweg das Core-Element über Lauf, Claim und Unit mit der
Quellstelle. Für dialogische und GitHub-Eingaben bilden dagegen
das jeweilige Eingangsartefakt und die dokumentierte Übernahme den
Zugang. Herkunftsangaben, Historie und archivierte Pläne ergänzen
damit die fachlichen Beziehungen aus Tabelle~\ref{t:core-objektmodell}.

Diese Angaben werden bei Fortschreibungen nicht überall einheitlich
nachgeführt. Eine Bestätigung kann Claim-Kennungen ergänzen, ohne
den eigenen Eingangs- und Laufbezug eindeutig am Element anzubinden.
Bei Verfeinerungen können aktueller Text und Herkunftsangaben
unterschiedliche Stände betreffen. Ein auffindbarer Ursprung erklärt
deshalb nicht automatisch jede spätere Ergänzung.
Abschnitt~\ref{sec:eval-provenienz} untersucht direkte Bezüge,
Beitragspfade und aktuelle Änderungsbelege getrennt.

\emph{Zustandsdimensionen.} Mehrere getrennte Angaben beschreiben
ein Element. Die \emph{Gültigkeit} unterscheidet geltende und
abgelöste Fassungen; Ablösung ist keine Löschung. Der
\emph{Fortschritt} bezeichnet bei Arbeitspaketen offene oder
erledigte Arbeit: Ein erledigtes Arbeitspaket kann weiterhin gültig
sein. Eine \emph{Blockade} kennzeichnet ungeklärte Voraussetzungen
und kann seine operative Weitergabe verhindern. Der
\emph{Bestätigungsmarker} unterscheidet etwa menschlich bestätigte
Elemente von solchen, die bei der Initialisierung verankert wurden.
Der \emph{Entscheidungsstand} hält fest, ob eine fachliche Frage
offen oder aufgelöst ist. Eine kompakte Statusanzeige fasst diese
Angaben zusammen, ersetzt ihre unterschiedlichen Bedeutungen aber nicht.

\emph{Fortschreibung und Auflösung.} Änderungen können ein
Element verfeinern oder durch eine neue Fassung ersetzen.
Widersprüche und unbeantwortete Fragen werden als
Entscheidungselemente fortgeführt. Abbildung~\ref{fig:entscheidungs-lifecycle} zeigt den
Fall, in dem eine spätere Auflösung fachliche Inhalte verändert und
abhängige Arbeit betreffen kann.

\begin{figure}[htbp]
  \centering
  \includegraphics[width=\textwidth]{figures/05/Abbildung5.5.pdf}
  \caption[Entscheidungsauflösung und abhängige Fortschreibung]{
    Entscheidungsauflösung mit fachlicher Änderung und möglicher
    abhängiger Fortschreibung (H: menschlicher Entscheidungspunkt;
    eigene Darstellung).}
  \label{fig:entscheidungs-lifecycle}
\end{figure}

Bei einem bestehenden Zielelement unterscheidet der Auflösungsplan
drei Ergebnisse: bisherige Aussage beibehalten, neue Aussage als
ablösende Fassung übernehmen oder bestehende Aussage verfeinern.
Er beschreibt das Ergebnis und die betroffenen Elemente; der Mensch
autorisiert die vorgeschlagenen Operationen am Entscheidungspunkt.
Die anschließende Anwendung erhält beim Verfeinern die Kennung und
erhöht die Version. Beim Ersatz entsteht dagegen ein neues Element
mit Ablösungsbeziehung; die frühere Fassung bleibt historisch erhalten.
Für freistehende Fragen ohne Zielelement gelten engere Regeln:
Sie können als geklärt geschlossen werden; bei entsprechend
gekennzeichneten Architekturfragen ist außerdem ein ausdrücklich
dokumentierter Verzicht auf eine fachliche Festlegung vorgesehen.
Ohne Zielelement gibt es nichts zu verfeinern oder abzulösen.

Abhängige Arbeit wird anhand der Beziehungen ermittelt. Nach einer
fachlichen Änderung markiert die Anwendung betroffene Arbeitspakete
mit Klärungsbedarf; der nachgelagerte Angleichungspfad greift diese
Markierung auf und legt seine Vorschläge erneut zur Entscheidung vor.
Eine Auflösung führt daher nicht automatisch zur abgeschlossenen
Angleichung aller abhängigen Inhalte. Der in
Abschnitt~\ref{sec:einordnende-fortschreibung} eingeordnete
Experimentfall E47-02 zeigt genau diese Grenze; seine Bewertung
steht in Abschnitt~\ref{sec:eval-fortschreibung}. Zusätzliche
Projektinformationen treten dagegen erneut über die reguläre
Einordnung ein; ihr Ergebnis ist nicht vorbestimmt.

Abgelöste Fassungen, aufgelöste Widerspruchsbeziehungen und
Ablehnungswissen bleiben erhalten. Frühere Ablehnungen können die
erneute Prüfung informieren, ersetzen aber keine neue menschliche
Entscheidung. Auch die modellgestützte Einordnung garantiert keine
fachlich richtige Klassifikation; Kapitel~\ref{chap:evaluation}
untersucht ausgewählte Fälle, keine allgemeine Fehlerhäufigkeit.
Welche Kontrollen vor der Übernahme einer Änderung gelten,
erklärt der folgende Abschnitt.
